% !TEX program = xelatex
% !TEX encoding = UTF-8
% ---------------------------------------------------------------
%  LaTeX 编译环境测试文件
%  编译方式（VSCode LaTeX Workshop 中任选一种）：
%    1) 快捷键 Ctrl+Alt+B 编译，Ctrl+Alt+V 预览
%    2) 左侧 TeX 面板 -> Build LaTeX project -> 选择 "xelatex" 配方
%  命令行：xelatex -synctex=1 -interaction=nonstopmode test.tex
% ---------------------------------------------------------------
\documentclass[UTF8, a4paper, 11pt, fontset=fandol]{ctexart}

% ---------- 页面与排版 ----------
\usepackage[top=2.5cm, bottom=2.5cm, left=2.5cm, right=2.5cm]{geometry}
\usepackage{indentfirst}          % 中文段落首行缩进
\setlength{\parindent}{2em}
\usepackage{setspace}
\onehalfspacing                    % 1.5 倍行距，论文常用

% ---------- 数学 ----------
\usepackage{amsmath, amssymb, bm}

% ---------- 图形表格 ----------
\usepackage{graphicx}
\usepackage{booktabs}
\usepackage{multirow}
\usepackage{caption}
\captionsetup{font=small}

% ---------- 其他 ----------
\usepackage{xcolor}
\usepackage{enumitem}
\usepackage{tikz}
\usetikzlibrary{arrows.meta, positioning, calc}
\usepackage{iftex}                 % 引擎检测
\usepackage[hidelinks, colorlinks=true, linkcolor=blue, citecolor=red]{hyperref}

% ---------- 自定义命令：显示当前编译引擎 ----------
\newcommand{\enginename}{%
  \ifXeTeX Xe\LaTeX%
  \else\ifPDFTeX pdf\LaTeX%
  \else\ifLuaTeX Lua\LaTeX%
  \else 未知引擎%
  \fi\fi\fi}

% ===============================================================
\title{\textbf{LaTeX 编译环境测试文档}\\[4pt]
       \large 中文 · 公式 · 表格 · 插图 · 交叉引用 · 参考文献}
\author{锦松}
\date{\today}

\begin{document}
\maketitle

\begin{abstract}
本文件用于验证 \TeX\ Live 编译链路是否正常工作，覆盖中文排版、数学公式、
表格、\textsc{TikZ} 插图、交叉引用与参考文献等常用功能。
若本文件能顺利生成 PDF，说明 VSCode + \LaTeX\ Workshop + \TeX\ Live 环境配置正确。
\end{abstract}

\tableofcontents
\bigskip

% ===============================================================
\section{编译环境自检}
\label{sec:env}

当前文档由 \textbf{\enginename} 引擎编译，编码为 UTF-8，
中文字体使用 \texttt{ctex} 宏包的 \texttt{fandol} 字体集（\TeX\ Live 自带，无需系统字体）。
若你更想使用 Windows 系统字体（如宋体、等线），把文档类选项改为
\texttt{fontset=windows} 即可。

表~\ref{tab:env} 列出了建议的编译配置。

\begin{table}[htbp]
  \centering
  \caption{建议的编译环境配置}
  \label{tab:env}
  \begin{tabular}{@{}ll@{}}
    \toprule
    项目     & 配置值 \\
    \midrule
    发行版   & \TeX\ Live 2026 \\
    编译引擎 & \enginename（推荐 Xe\LaTeX） \\
    文件编码 & UTF-8（VSCode 右下角确认） \\
    文档类   & \texttt{ctexart} / \texttt{ctexbook} \\
    中文支持 & \texttt{ctex} 宏包 + \texttt{fontset=fandol} \\
    参考文献 & \texttt{biblatex} 或 \texttt{thebibliography} \\
    \bottomrule
  \end{tabular}
\end{table}

% ===============================================================
\section{中文与英文混排测试}

中文测试：知识投毒攻击通过向检索语料中注入对抗性文本，
诱导大语言模型生成攻击者预设的答案。英文混排测试：
\textit{Retrieval-Augmented Generation (RAG) improves factuality by grounding
generation in external corpora}，但同时扩大了攻击面。

中英文之间自动添加的间距应由 \texttt{ctex} 处理，无需手动输入空格。
数字与单位示例：检索 Top-$k = 5$，投毒比例 $0.08\%$，ASR 从 $0.833$ 降到 $0.217$。

% ===============================================================
\section{数学公式测试}

行内公式示例：检索增强生成的输出可写为
$p_\theta(y \mid q, \mathcal{D}_k)$，其中 $\mathcal{D}_k = \{d_1, \dots, d_k\}$
为检索到的 top-$k$ 文档集合。

带编号的行间公式（攻击成功率定义）：
\begin{equation}
  \mathrm{ASR} = \frac{1}{|\mathcal{Q}|} \sum_{q \in \mathcal{Q}}
  \mathbb{1}\!\left[\, \hat{y}_q \in \mathcal{A}_q \,\right],
  \label{eq:asr}
\end{equation}
其中 $\hat{y}_q$ 为模型对查询 $q$ 的生成答案，$\mathcal{A}_q$ 为攻击者目标答案集合，
$\mathbb{1}[\cdot]$ 为指示函数。式~(\ref{eq:asr}) 是全文评价的核心指标。

多行对齐环境：
\begin{align}
  \mathrm{CCI}(c) &= \frac{1}{|c|}\sum_{d \in c} w_d \cdot \mathrm{JS}(d \parallel r),
  \label{eq:cci1} \\[4pt]
  \mathrm{score}(c) &= \alpha \cdot \mathrm{CCI}(c) - \beta \cdot \mathrm{Supp}(c),
  \label{eq:cci2}
\end{align}
其中 $\mathrm{JS}(\cdot \parallel \cdot)$ 为 Jensen--Shannon 散度，
$\mathrm{Supp}(c)$ 表示证据联盟 $c$ 的独立支撑度，
$\alpha, \beta > 0$ 为权衡系数。

矩阵与希腊字母：
\[
  \bm{W} =
  \begin{pmatrix}
    w_{11} & w_{12} & \cdots & w_{1n} \\
    w_{21} & w_{22} & \cdots & w_{2n} \\
    \vdots & \vdots & \ddots & \vdots \\
    w_{m1} & w_{m2} & \cdots & w_{mn}
  \end{pmatrix},
  \qquad
  \theta^{*} = \arg\max_{\theta}\ \mathcal{L}(\theta; \lambda, \gamma, \Omega).
\]

% ===============================================================
\section{表格测试}

表~\ref{tab:result} 是一个典型的实验结果表格样式（数字仅为占位示例，
\textbf{不可用于论文正文}）。

\begin{table}[htbp]
  \centering
  \caption{不同防御方法的攻击成功率与拒答率对比（示例数据）}
  \label{tab:result}
  \begin{tabular}{@{}lcc@{}}
    \toprule
    方法 & ASR $\downarrow$ & 拒答率 $\downarrow$ \\
    \midrule
    无防御 (B0)        & 0.833 & 0.017 \\
    相关性过滤 (B1)    & 0.833 & 0.017 \\
    逐文档 LOO (B6)    & 0.650 & 0.067 \\
    \textbf{本文方法}  & \textbf{0.217} & 0.400 \\
    \bottomrule
  \end{tabular}
\end{table}

跨列跨行示例见表~\ref{tab:merge}。

\begin{table}[htbp]
  \centering
  \caption{跨行跨列示例}
  \label{tab:merge}
  \begin{tabular}{@{}lccc@{}}
    \toprule
    \multirow{2}{*}{数据集} & \multicolumn{2}{c}{ASR} & \multirow{2}{*}{平均} \\
    \cmidrule(lr){2-3}
                       & Qwen2.5-7B & Llama-3.1-8B & \\
    \midrule
    NQ       & 0.833 & 0.650 & 0.742 \\
    HotpotQA & 0.850 & 0.667 & 0.759 \\
    MS~MARCO & 0.817 & 0.633 & 0.725 \\
    \bottomrule
  \end{tabular}
\end{table}

% ===============================================================
\section{插图测试}

\subsection{TikZ 绘图（无需外部图片文件）}

图~\ref{fig:curve} 用 \textsc{TikZ} 直接绘制，不依赖任何外部图片，
因此一定能编译通过。

\begin{figure}[htbp]
  \centering
  \begin{tikzpicture}[scale=0.95]
    % 坐标轴
    \draw[->, thick] (0,0) -- (6.6,0) node[right] {联盟规模 $|c|$};
    \draw[->, thick] (0,0) -- (0,3.8) node[above] {影响力};
    % 曲线
    \draw[domain=0.2:6.2, samples=120, thick, blue!70!black]
         plot (\x, {2.6 * exp(-(\x - 5.2)^2 / 3.2) + 0.35});
    % 标记点
    \foreach \x/\y in {1/0.6, 2/0.95, 3/1.55, 4/2.25, 5/2.85} {
      \filldraw[red!80!black] (\x,\y) circle (2.2pt);
    }
    % 刻度
    \foreach \x in {1,2,3,4,5} {
      \draw (\x,0) -- (\x,-0.1) node[below] {\small \x};
    }
    \foreach \y in {1,2,3} {
      \draw (0,\y) -- (-0.1,\y) node[left] {\small \y};
    }
    \node[above right, blue!70!black] at (3.4,2.9) {\small 影响力随联盟规模增长};
  \end{tikzpicture}
  \caption{TikZ 绘制的示例曲线}
  \label{fig:curve}
\end{figure}

图~\ref{fig:flow} 是一个流程框图，验证 \texttt{positioning} 与箭头库是否正常。

\begin{figure}[htbp]
  \centering
  \begin{tikzpicture}[
      node distance=1.5cm,
      box/.style={draw, rounded corners=3pt, fill=blue!8,
                  minimum width=2.1cm, minimum height=0.85cm, font=\small},
      arw/.style={-{Stealth[length=2.5mm]}, thick}
    ]
    \node[box] (q)  {查询 $q$};
    \node[box, right=of q] (r) {检索 Top-$k$};
    \node[box, right=of r] (g) {联盟构建};
    \node[box, below=of g] (f) {反事实过滤};
    \node[box, left=of f]  (o) {生成答案};
    \draw[arw] (q) -- (r);
    \draw[arw] (r) -- (g);
    \draw[arw] (g) -- (f);
    \draw[arw] (f) -- (o);
    \draw[arw, dashed, red!70] (o.west) -- ++(-1.2,0) |- (q.west)
          node[pos=0.25, above, font=\scriptsize, black] {反馈};
  \end{tikzpicture}
  \caption{TikZ 绘制的流程框图}
  \label{fig:flow}
\end{figure}

\subsection{外部图片（可选验证）}

如果你有图片文件，取消下面这段注释以验证 \texttt{graphicx} 是否正常
（把文件名换成真实路径）：

\begin{verbatim}
\begin{figure}[htbp]
  \centering
  \includegraphics[width=0.6\textwidth]{figures/demo.png}
  \caption{外部图片示例}
  \label{fig:external}
\end{figure}
\end{verbatim}

% ===============================================================
\section{列表与交叉引用测试}

无序列表：
\begin{itemize}[noitemsep]
  \item 中文条目一：知识投毒
  \item 中文条目二：证据联盟
  \item 英文条目：\texttt{coalition-level attribution}
\end{itemize}

有序列表：
\begin{enumerate}[label=\arabic*), noitemsep]
  \item 检索候选文档并执行关系抽取；
  \item 构建证据联盟并计算控制力指标；
  \item 对高控制力、低支撑度的联盟执行反事实过滤；
  \item 用过滤后的证据集生成最终答案。
\end{enumerate}

交叉引用检查：\autoref{sec:env} 节、表~\ref{tab:result}、式~(\ref{eq:asr})、
图~\ref{fig:curve} 与图~\ref{fig:flow}。如果这些都显示为正确编号（而非 \texttt{??}），
说明需要编译两次的流程已正常完成。

% ===============================================================
\section{参考文献测试}

本例使用 \texttt{thebibliography} 环境，不依赖 Bib\TeX，可单次编译出结果。
引用示例：PoisonedRAG~\cite{zou2025poisonedrag} 揭示了 RAG 的投毒脆弱性；
TrustRAG~\cite{xiang2024trustrag} 采用聚类过滤思路；
反事实推断相关方法见~\cite{pearl2009causality}。

\begin{thebibliography}{9}
  \bibitem{zou2025poisonedrag}
  W.~Zou, R.~Geng, B.~Wang, et al.
  \newblock PoisonedRAG: Knowledge Corruption Attacks to Retrieval-Augmented
           Generation of Large Language Models.
  \newblock \emph{USENIX Security Symposium}, 2025.

  \bibitem{xiang2024trustrag}
  C.~Xiang, T.~Wu, et al.
  \newblock Certifiably Robust RAG against Retrieval Corruption.
  \newblock \emph{arXiv preprint}, 2024.

  \bibitem{pearl2009causality}
  J.~Pearl.
  \newblock \emph{Causality: Models, Reasoning and Inference}.
  \newblock Cambridge University Press, 2nd edition, 2009.
\end{thebibliography}

% ===============================================================
\section*{编译状态总结}

\begin{itemize}[noitemsep]
  \item 若 PDF 正常生成、中文无乱码、公式与图表均显示正确，则环境配置成功。
  \item 若出现 \texttt{??} 引用，再编译一次即可（交叉引用需要两轮）。
  \item 若中文显示为空白或方框，检查是否使用了 Xe\LaTeX\ 引擎，
        或把 \texttt{fontset=fandol} 改为 \texttt{fontset=windows}。
\end{itemize}

\end{document}
